% Part: sets-functions-relations
% Chapter: infinite
% Section: dedekind-induction

\documentclass[../../../include/open-logic-section]{subfiles}

\begin{document}

\olfileid{sfr}{infinite}{induction}
\olsection[গাণিতিক আরোহ]{ডেডেকিন্ড বীজগঠন ও গাণিতিক আরোহ}

এখন মূল কথাটি হলো, একটি ডেডেকিন্ড বীজগঠন—বস্তুত, \emph{যেকোনো} ডেডেকিন্ড বীজগঠন—স্বাভাবিক সংখ্যার প্রতিস্থাপক হিসেবে কাজ করবে। নিচের সরল অনুসিদ্ধান্তটির কারণেই তা সম্ভব:

\begin{thm}[গাণিতিক আরোহ]\ollabel{thm:dedinfiniteinduction}
ধরি $N, s, o$ মিলে একটি ডেডেকিন্ড বীজগঠন গঠন করে। তাহলে যেকোনো সেট $X$-এর জন্য:
	\begin{center}
		যদি $o \in X$ এবং $(\forall {n} \in N \cap X){s}({n}) \in X$ হয়, {তবে} $N \subseteq X$।
	\end{center}
\end{thm}

\begin{proof}
ডেডেকিন্ড বীজগঠনের সংজ্ঞা অনুসারে $N = \closureofunder{s}{o}$। এবার ${o} \in X$ এবং $(\forall {n} \in N)(n \in X \lif
{s}({n}) \in X)$—দুটিই সত্য হলে বাহক সেট ও প্রদত্ত সেটের ছেদে নির্দিষ্ট আদিবিন্দুটি আছে এবং ছেদটি উত্তরসূরি অপেক্ষকের অধীনে বদ্ধ: ছেদটির যেকোনো উপাদানের উত্তরসূরি বাহক সেটে থাকে, কারণ সেটটি নিজেই উত্তরসূরি অপেক্ষকের অধীনে বদ্ধ, আর প্রদত্ত শর্তে সেটি অন্য সেটেও থাকে। তাই আবরণের ক্ষুদ্রতমতার ধর্মে বাহক সেটটি ওই ছেদের উপসেট; বিশেষত $N = \closureofunder{s}{o} \subseteq X$।
\end{proof}
% BN-SRC-017 TARGET-CORRECTION-BEGIN
\par\smallskip\noindent\textnormal{\small\textbf{উৎস-সংশোধন (BN-SRC-017):} হিমায়িত ইংরেজি উৎসটি প্রদত্ত শর্ত থেকে সরাসরি $\closureofunder{s}{o}\subseteq X$ বলে, কিন্তু $X$-এর বাইরের উপাদানের উপর $s$-এর আচরণ নিয়ন্ত্রিত নয় বলে $X$ নিজে $s$-বদ্ধ হওয়া অনুসৃত হয় না। বাংলা প্রমাণে $N\cap X$ ব্যবহার করা হয়েছে: সেটি $o$-কে ধারণ করে এবং $N$ ও প্রদত্ত শর্ত একত্রে সেটিকে $s$-বদ্ধ করে; তাই ক্ষুদ্রতমতার ধর্ম প্রযোজ্য।} % BN-SRC-017 TARGET-CORRECTION-END

আরোহ স্বাভাবিক সংখ্যার বৈশিষ্ট্যসূচক বলে কথাটি দাঁড়ায় এই: যেকোনো ডেডেকিন্ড-অসীম সেট থেকে আমরা একটি ডেডেকিন্ড বীজগঠন তৈরি করতে পারি এবং সেই বীজগঠনকে স্বাভাবিক সংখ্যার প্রতিস্থাপক হিসেবে ব্যবহার করতে পারি।

স্বীকার করতেই হয়, \olref{thm:dedinfiniteinduction}-এ আরোহকে \emph{সেটতাত্ত্বিক} ভাষায় সূত্রবদ্ধ করা হয়েছে। কিন্তু নীতিটিকে আরও পরিচিত ভাষায় সহজেই লেখা যায়:

\begin{cor}\ollabel{natinductionschema}
ধরি $N, s, o$ মিলে একটি ডেডেকিন্ড বীজগঠন গঠন করে। তাহলে পরামিতিযুক্ত হতে পারে এমন যেকোনো সূত্র $\phi(x)$-এর জন্য:
\begin{center}
	যদি $\phi(o)$ এবং $(\forall {n} \in N)(\phi(n)\lif
	\phi({s}({n})))$ হয়, {তবে} $(\forall n \in N)\phi(n)$
\end{center}
\end{cor}

\begin{proof}
ধরি $X = \Setabs{n \in N}{\phi(n)}$; এবার
\olref{thm:dedinfiniteinduction} ব্যবহার করো।
\end{proof}

এই ফলে আমরা একটি সূত্রের ``পরামিতি থাকা'' নিয়ে কথা বলেছি। মোটামুটিভাবে এর অর্থ হলো, যেকোনো বস্তু $c_1, \ldots, c_k$-এর জন্য আমরা $\phi(x, c_1, \ldots, c_k)$ নিয়ে কাজ করতে পারি। আরও নির্ভুলভাবে, ``পরামিতি'' উল্লেখ না করেই ফলটি এভাবে বলা যায়। যেকোনো সূত্র $\phi(x, v_1, \ldots, v_k)$, যার সব মুক্ত চলরাশি দেখানো হয়েছে, তার জন্য পাই:
	\begin{align*}
			\forall v_1 \ldots \forall v_k((&\phi(o, v_1,\ldots, v_k) \land {}\\
			&	(\forall x \in N)(\phi(x,v_1, \ldots, v_k) \lif \phi(s(x), v_1,\ldots, v_k))) \lif {}\\
			&\hspace{3em} (\forall x \in N)\phi(x, v_1,\ldots, v_k))
	\end{align*}
স্পষ্টতই, ``পরামিতি থাকা'' কথাটি পাঠকে অনেক সহজ করে তুলতে পারে। (\olref[sth][][]{part}-এ এই রীতিটি আমরা প্রায়ই ব্যবহার করব।)

ডেডেকিন্ড বীজগঠনে ফিরে আসি: যেকোনো ডেডেকিন্ড বীজগঠন দেওয়া থাকলে যোগ, গুণ ও ঘাতের প্রচলিত গাণিতিক অপেক্ষকগুলিও আমরা সংজ্ঞায়িত করতে পারি। তবে কাজটি সহজ নয়; এতে \emph{পুনরাবৃত্ত সংজ্ঞা}-র কৌশল লাগে। এই কৌশলটি আমরা অনেক পরে, আরও সাধারণ পরিপ্রেক্ষিতে প্রবর্তন ও ন্যায্যতা প্রতিপাদন করব। (আগ্রহী পাঠক \olref[sth][ord-arithmetic][]{chap}-এর পরে এই অংশে ফিরে আসতে পারেন, অথবা \citealt[pp.~95--8]{Potter2004}-এর মতো কোনো বিকল্প আলোচনা পড়তে পারেন।) তবে $N, s, o$ মিলে একটি ডেডেকিন্ড বীজগঠন গঠন করলে শেষ পর্যন্ত আমরা নিচের শর্তগুলি নির্ধারণ করতে পারব:
\begin{align*}
	{a} + {o} &= {a} & & & {a} \times {o} &= {o} & & & {a}^{o} &= s(o)\\
{a} + {s}({b}) &= {s}({a}+{b}) &&& {a} \times {s}({b}) &= ({a}\times {b}) + {a}  & & & {a}^{{s}({b})} &= {a}^{b} \times {a}
\end{align*}
এবং দেখাতে পারব যে এগুলি প্রত্যাশিতভাবেই আচরণ করে।

\end{document}
